\begin{definition}[Portability~\cite{sedova_high-performance_2018}]
An application can be said to be portable if the cost of porting is less than the cost of re-writing:
$$DP = 1 - \left(\frac{C_P}{C_R}\right)$$
where $C_P$ is the cost to port and $C_R$ the cost to rewrite the program. 
\end{definition}

\begin{itemize}
    \item $1$ $\rightarrow$ completely portable
    \item positive $\rightarrow$ porting is more profitable
    \item binary vs source code portability
\end{itemize}

\begin{table}[]
    \newcommand{\supported}{\textcolor{ForestGreen}{\faCheckCircle}}
    \newcommand{\partialsupported}{\textcolor{BurntOrange}{\faCheckCircle}}
    \newcommand{\unsupported}{\textcolor{Red}{\faTimesCircle}}
    \centering
    \begin{tblr}{colspec = {Q[l,m] *{6}{X[c,m]}},
                 row{1,2} = {font=\bfseries, bg=gray!50},
                 row{odd[2]} = {bg=gray!25},
                 rowsep=4pt,
                 }
        \toprule
                      & \SetCell[c=3]{c}\glspl{cpu} &&             & \SetCell[c=3]{c}\glspl{gpu} &&                            \\
                      & x86\_64      & POWER9       & ARM          & NVIDIA            & AMD               & Intel             \\\midrule
        OpenMP        & \supported   & \supported   & \supported   & \supported        & \supported        & \supported        \\
        CUDA          & \unsupported & \unsupported & \unsupported & \supported        & \unsupported      & \unsupported      \\
        HPX           & \supported   & \supported   & \supported   & \partialsupported & \partialsupported & \partialsupported \\
        OpenCL        & \supported   & \supported   & \supported   & \supported        & \supported        & \supported        \\
        Standard \cpp & \supported   & \supported   & \supported   & \supported        & \supported        & \supported        \\
        SYCL          & \supported   & \supported   & \supported   & \supported        & \supported        & \supported        \\
        Kokkos        & \supported   & \supported   & \supported   & \supported        & \supported        & \supported        \\
        HIP           & \unsupported & \unsupported & \unsupported & \supported        & \supported        & \unsupported      \\
        \bottomrule
    \end{tblr}
    \caption{%
    Either we can use two categories, simply splitting the hardware platforms in \glspl{cpu} and \glspl{gpu}, or split it into more fine-grain categories like different \gls{cpu} and \gls{gpu} architectures. 
    Additional, other hardware platforms inside these categories, e.g., RISC-V \glspl{cpu} or ARM \glspl{gpu} or even completely new categories like \glspl{fpga} or \glspl{tpu} are also possible. 
    For \gls{hpx}, \glspl{gpu} are supported via executors, but the actual compute kernels must be written using another framework like \gls{cuda}, \gls{hip}, SYCL, or Kokkos.
    Note that for the standardized frameworks, it can happen that no single implementation supports all hardware platforms, but only a mixture does. 
    }
    \label{tab:portability_categories}
\end{table}



\begin{table}[]
    \newcommand{\supported}{\textcolor{ForestGreen}{\faCheckCircle}}
    \newcommand{\partialsupported}{\textcolor{BurntOrange}{\faCheckCircle}}
    \newcommand{\unsupported}{\textcolor{Red}{\faTimesCircle}}
    \centering
    \begin{tblr}{colspec = {Q[l,m] *{4}{X[c,m]}},
                 row{1,2} = {font=\bfseries, bg=gray!50},
                 row{odd[2]} = {bg=gray!25},
                 rowsep=4pt,
                 cell{1}{1}={r=2}{l},
                 cell{1}{2}={r=2}{c},
                 }
        \toprule
        {programming\\framework} & \glspl{cpu}  & \SetCell[c=3]{c}\glspl{gpu} &&                            \\
                                 &              & NVIDIA            & AMD               & Intel             \\\midrule
        OpenMP (target offload)  & \supported   & \supported        & \supported        & \supported        \\
        CUDA                     & \unsupported & \supported        & \unsupported      & \unsupported      \\
        HPX                      & \supported   & \partialsupported & \partialsupported & \partialsupported \\
        OpenCL                   & \supported   & \supported        & \supported        & \supported        \\
        Standard \cpp            & \supported   & \supported        & \supported        & \supported        \\
        SYCL                     & \supported   & \supported        & \supported        & \supported        \\
        Kokkos                   & \supported   & \supported        & \supported        & \supported        \\
        HIP                      & \unsupported & \supported        & \supported        & \unsupported      \\
        \bottomrule
    \end{tblr}
    \caption{%
    The list of the programming frameworks used in the remainder of this dissertation including their supported target platforms. 
    Five out of the investigated eight programming frameworks promise to support all major hardware platforms, i.e., \glspl{cpu} and \glspl{gpu} from NVIDIA, \gls{amd}, and Intel. 
    Note that we do not make use of \gls{openmp}'s target offload feature and, therefore, later only investigate \gls{openmp} on \glspl{cpu}. 
    For \gls{hpx}, \glspl{gpu} are supported via executors, but the actual compute kernels must be written using another framework like \gls{cuda}, \gls{hip}, SYCL, or Kokkos.
    % Note that for the standardized frameworks, it can happen that no single implementation supports all hardware platforms, but only a mixture does. 
    }
    \label{tab:frameworks_overview}
\end{table}